% ============================================================
% NUEVA SUBSECCION: Metricas ampliadas (7 metricas)
% Pegar dentro de la Seccion VII (Resultados), despues de las Tablas I-II
% originales, o como nueva subseccion "D. Metricas ampliadas".
% Requiere \usepackage{graphicx} y \usepackage{booktabs} (ya incluidos en main.tex)
% ============================================================

\subsection{M\'etricas ampliadas: memoria y crecimiento emp\'irico}

Con el fin de caracterizar con mayor profundidad el comportamiento de BMSSP
frente a Dijkstra, se repiti\'o la experimentaci\'on del Experimento 1 y el
Experimento 2 con \textbf{10 repeticiones por instancia} (en lugar de una sola
ejecuci\'on), y se instrumentaron dos m\'etricas adicionales de memoria: la
memoria pico del proceso (\emph{peak RSS}, v\'ia \texttt{getrusage()}) y los
bytes de heap asignados durante la fase de soluci\'on (v\'ia una sobrecarga de
\texttt{operator new}/\texttt{delete} activa solo durante la llamada al
algoritmo). Cada algoritmo se ejecut\'o en un proceso independiente para evitar
contaminaci\'on de memoria entre mediciones.

Con 10 repeticiones, el speedup de BMSSP respecto a Dijkstra convergi\'o a un
rango de \textbf{3.47x--4.13x} en las instancias m\'as grandes, sensiblemente
m\'as estable que el rango 3.1x--10.2x obtenido con una \'unica ejecuci\'on, y
consistente con el factor de 3--4x reportado por Castro et al.
\cite{implbmssp2025}.

En cuanto a memoria, se observa una disociaci\'on interesante entre ambas
m\'etricas: la memoria pico (RSS) de BMSSP super\'o a la de Dijkstra en solo un
7.5\% en la instancia m\'as grande de Exp2 ($m=250\,000$), mientras que el
volumen de bytes de heap asignados durante la soluci\'on fue, en promedio,
\textbf{23 veces mayor} en BMSSP. Esto sugiere que la sobrecarga de BMSSP no
proviene de un mayor footprint sostenido, sino de un volumen alto de
asignaciones y liberaciones transitorias asociadas a su estructura recursiva
por niveles (cola de lotes, mapas de pivotes), consistente con el dise\~no del
algoritmo descrito en la Secci\'on III-A2.

Un ajuste de ley de potencia ($t = c \cdot n^p$) sobre los tiempos medios del
Experimento 1 produjo:
\begin{align}
t_{\text{Dijkstra}} &\approx 0.18 \cdot n^{1.12} \quad (R^2 = 0.997) \\
t_{\text{BMSSP}} &\approx 5.48 \cdot n^{0.86} \quad (R^2 = 0.973)
\end{align}
Extrapolando ambos modelos, se estima un punto de cruce te\'orico en
$n \approx 5.3 \times 10^5$. Sin embargo, esta cifra debe interpretarse con
cautela: proviene de una extrapolaci\'on de un modelo emp\'irico ajustado sobre
un rango de solo $n\in[100,6400]$, casi dos \'ordenes de magnitud por debajo
del punto proyectado, y difiere sustancialmente del estimado por Castro et al.
\cite{implbmssp2025} ($n\approx10^{67}$), obtenido mediante an\'alisis
asint\'otico rigoroso de las constantes del algoritmo en lugar de extrapolaci\'on
emp\'irica. Ambas cifras coinciden, no obstante, en la conclusi\'on cualitativa:
BMSSP no superar\'ia a Dijkstra dentro de ning\'un rango de tama\~no
computacionalmente alcanzable en la pr\'actica.

\begin{table}[t]
\caption{M\'etricas ampliadas, Experimento 1 (10 repeticiones)}
\label{tab:metricas7}
\centering
\small
\begin{tabular}{@{}rrrrr@{}}
\toprule
$n$ & Speedup & Heap Dij. (KB) & Heap BMSSP (KB) & Raz\'on \\
\midrule
100    & 9.44x & 8.8   & 291.5   & 33.3x \\
200    & 9.49x & 17.5  & 601.7   & 34.3x \\
400    & 5.45x & 35.1  & 949.0   & 27.0x \\
800    & 4.39x & 70.2  & 1704.2  & 24.3x \\
1\,600 & 4.49x & 140.5 & 3259.1  & 23.2x \\
3\,200 & 3.82x & 281.0 & 6311.0  & 22.5x \\
6\,400 & 3.47x & 562.0 & 12035.0 & 21.4x \\
\bottomrule
\end{tabular}
\end{table}

% Para insertar las figuras (copiar los .png a la carpeta figs/ del proyecto):
% \begin{figure}[t]
% \centering
% \includegraphics[width=0.95\linewidth]{figs/fig_memoria_estructura_exp1.png}
% \caption{Bytes de heap asignados durante la fase de soluci\'on (Experimento 1).}
% \label{fig:heap}
% \end{figure}
